\chapter{The Solution of Differential Equations}

\section{Initial-Value Problems}

\begin{definition}[Ordinary differential equation]
An ordinary differential equation (ODE) is an equation involving an unknown function of one independent variable and one or more of its derivatives.
\end{definition}

An initial-value problem (IVP) specifies the value of the solution, and possibly some derivatives, at one point.  A boundary-value problem specifies conditions at different points.

\begin{definition}[First-order IVP]
A solution of
\[
 y'=f(x,y),
 \qquad y(x_0)=y_0,
\]
on an interval $[a,b]$ is a differentiable function $y=\phi(x)$ such that
\[
 \phi(x_0)=y_0,
 \qquad
 \phi'(x)=f(x,\phi(x))
\]
for every $x\in[a,b]$.
\end{definition}

\begin{definition}[Lipschitz condition]
A function $f(x,y)$ satisfies a Lipschitz condition in $y$ on a set $D$ if a constant $L>0$ exists such that
\[
 |f(x,y_1)-f(x,y_2)|\le L|y_1-y_2|
\]
whenever $(x,y_1),(x,y_2)\in D$.
\end{definition}

\begin{theorem}[Existence and uniqueness]
Suppose $f$ is continuous on
\[
 D=\{(x,y):a\le x\le b,\ -\infty<y<\infty\}
\]
and satisfies a Lipschitz condition in $y$.  Then the IVP
\[
 y'=f(x,y),
 \qquad y(a)=y_0,
\]
has a unique solution on $[a,b]$.
\end{theorem}

A convenient sufficient condition is a bounded partial derivative:
if $D$ is convex and
\[
 \left|\frac{\partial f}{\partial y}(x,y)\right|\le L
\]
on $D$, then $f$ is Lipschitz in $y$ with constant $L$.

\begin{example}
Show that the IVP
\[
 y'=1+x\sin(xy),
 \qquad y(0)=0,
 \qquad0\le x\le2,
\]
has a unique solution.
\end{example}

\begin{solution}[3.8cm]
The function
\[
 f(x,y)=1+x\sin(xy)
\]
is continuous.  Also
\[
 \frac{\partial f}{\partial y}=x^2\cos(xy),
\]
so on $0\le x\le2$,
\[
 \left|\frac{\partial f}{\partial y}\right|
 \le x^2\le4.
\]
Thus $f$ satisfies a Lipschitz condition in $y$ with $L=4$.  The existence-and-uniqueness theorem applies.
\end{solution}

A \emph{one-step method} uses only information from the current point $(x_n,y_n)$ to find $y_{n+1}$.  A \emph{multistep method} uses data from several previous points.

\section{Euler's Method}

For mesh points
\[
 x_n=x_0+nh,
\]
Euler's method is
\[
 y_{n+1}=y_n+hf(x_n,y_n).
\]
It follows from truncating the Taylor expansion
\[
 y(x_n+h)=y(x_n)+hy'(x_n)+\frac{h^2}{2}y''(\xi_n).
\]
The local truncation error is $\bigO(h^2)$ and the accumulated global error is $\bigO(h)$.

\begin{example}
Consider
\[
 y'=x+y,
 \qquad y(0)=1.
\]
\begin{enumerate}[label=(\alph*),leftmargin=2.7em]
  \item Use Euler's method with $h=0.1$ to approximate $y(0.5)$.
  \item Estimate the truncation error in the first step using the next two Taylor terms.
  \item If the exact values are $y(0.1)=1.11034$ and $y(0.5)=1.79744$, find the errors at those points.
\end{enumerate}
\end{example}

\begin{solution}[9.5cm]
The Euler recurrence is
\[
 y_{n+1}=y_n+0.1(x_n+y_n).
\]
It gives
\[
\begin{array}{c|c|c}
\toprule
n&x_n&y_n\\
\midrule
0&0.0&1.00000\\
1&0.1&1.10000\\
2&0.2&1.22000\\
3&0.3&1.36200\\
4&0.4&1.52820\\
5&0.5&1.72102\\
\bottomrule
\end{array}
\]
Hence $y(0.5)\approx1.72102$.

Since $y'=x+y$,
\[
 y''=1+y'=1+x+y,
 \qquad
 y'''=1+y'=1+x+y.
\]
At $x=0$, $y=1$, so $y''(0)=2$ and $y'''(0)=2$.  The next two Taylor terms are
\[
 \frac{h^2}{2}y''(0)+\frac{h^3}{6}y'''(0)
 =\frac{0.1^2}{2}(2)+\frac{0.1^3}{6}(2)
 =0.01033.
\]
The actual first-step error is
\[
 1.11034-1.10000=0.01034.
\]
At $x=0.5$, the absolute error is
\[
 |1.79744-1.72102|=0.07642.
\]
\end{solution}

\begin{example}
For the same IVP, compare Euler approximations using $h=0.1$ and $h=0.05$ with the exact solution
\[
 y(x)=2e^x-x-1.
\]
\end{example}

\begin{solution}[5.5cm]
The smaller step gives consistently smaller errors:
\[
\begin{array}{c|cc|c|cc}
\toprule
x&y_{h=0.1}&y_{h=0.05}&y(x)&E_{0.1}&E_{0.05}\\
\midrule
0.1&1.10000&1.10500&1.11034&0.01034&0.00534\\
0.2&1.22000&1.23101&1.24281&0.02281&0.01180\\
0.3&1.36200&1.38019&1.39972&0.03772&0.01953\\
0.4&1.52820&1.55491&1.58365&0.05545&0.02874\\
0.5&1.72102&1.75779&1.79744&0.07642&0.03965\\
\bottomrule
\end{array}
\]
Halving $h$ approximately halves the global error, consistent with Euler's first-order accuracy.
\end{solution}

\section{Improved Euler Method}

Improved Euler, also called Heun's method or the explicit trapezoidal method, uses a predictor and corrector:
\[
 y_{n+1}^{*}=y_n+hf(x_n,y_n),
\]
\[
 y_{n+1}=y_n+\frac h2
 \left[f(x_n,y_n)+f(x_{n+1},y_{n+1}^{*})\right].
\]
Its local error is $\bigO(h^3)$ and its global error is $\bigO(h^2)$.

\begin{example}
Use improved Euler with $h=0.1$ for
\[
 y'=x+y,
 \qquad y(0)=1,
\]
on $[0,1]$.  Compare the approximations with $y(x)=2e^x-x-1$.
\end{example}

\begin{solution}[10.0cm]
The results are
\[
\begin{array}{c|c|c|c}
\toprule
x_k&y_k&y(x_k)&|y(x_k)-y_k|\\
\midrule
0.0&1.00000&1.00000&0.00000\\
0.1&1.11000&1.11034&0.00034\\
0.2&1.24205&1.24281&0.00076\\
0.3&1.39847&1.39972&0.00125\\
0.4&1.58180&1.58365&0.00184\\
0.5&1.79489&1.79744&0.00254\\
0.6&2.04086&2.04424&0.00338\\
0.7&2.32315&2.32751&0.00436\\
0.8&2.64558&2.65108&0.00550\\
0.9&3.01237&3.01921&0.00684\\
1.0&3.42817&3.43656&0.00839\\
\bottomrule
\end{array}
\]
The errors are much smaller than those from Euler's method with the same step size because the method is second order globally.
\end{solution}

\section{Taylor Series Methods}

A Taylor method of order $p$ uses
\[
 y_{n+1}=y_n+hy'_n+\frac{h^2}{2!}y''_n+\frac{h^3}{3!}y^{(3)}_n+
 \cdots+\frac{h^p}{p!}y_n^{(p)}.
\]
Its local truncation error is $\bigO(h^{p+1})$ and its global truncation error is $\bigO(h^p)$.  The main drawback is that higher derivatives must be derived and evaluated.

\begin{example}
For
\[
 y'=2xy,
 \qquad y(1)=1,
\]
use a second-order Taylor method with $h=0.1$ to approximate $y(1.2)$.
\end{example}

\begin{solution}[5.0cm]
Since
\[
 y'=2xy,
\]
we have
\[
 y''=2y+2xy'=2y+4x^2y=2(1+2x^2)y.
\]
The recurrence is
\[
 y_{n+1}=y_n+0.1(2x_ny_n)
 +\frac{0.1^2}{2}\,2(1+2x_n^2)y_n.
\]
At $(x_0,y_0)=(1,1)$,
\[
 y_1=1+0.2+0.03=1.23.
\]
At $(x_1,y_1)=(1.1,1.23)$,
\[
 y'_1=2.706,
 \qquad y''_1=8.4132,
\]
so
\[
 y_2=1.23+0.1(2.706)+0.005(8.4132)
 =1.542666.
\]
Therefore
\[
 y(1.2)\approx1.5427.
\]
\end{solution}

\section{Runge--Kutta Methods}

Runge--Kutta methods achieve high order without explicitly differentiating $f$.  The classical fourth-order method is
\[
 y_{n+1}=y_n+\frac16(k_1+2k_2+2k_3+k_4),
\]
where
\[
 k_1=hf(x_n,y_n),
\]
\[
 k_2=hf\left(x_n+\frac h2,y_n+\frac{k_1}{2}\right),
\]
\[
 k_3=hf\left(x_n+\frac h2,y_n+\frac{k_2}{2}\right),
\]
\[
 k_4=hf(x_n+h,y_n+k_3).
\]
RK4 has local error $\bigO(h^5)$ and global error $\bigO(h^4)$.

\begin{example}
Use RK4 with $h=0.1$ to approximate $y(0.2)$ for
\[
 y'=x+y,
 \qquad y(0)=1.
\]
\end{example}

\begin{solution}[6.2cm]
For the first step,
\[
 k_1=0.100000,
 \quad k_2=0.110000,
 \quad k_3=0.110500,
 \quad k_4=0.121050,
\]
which gives
\[
 y_1=1.110341667.
\]
For the second step,
\[
 k_1=0.121034167,
 \quad k_2=0.132085875,
 \quad k_3=0.132638460,
 \quad k_4=0.144298013.
\]
Hence
\[
 y_2=1.242805142.
\]
Therefore
\[
 y(0.2)\approx1.2428.
\]
\end{solution}

\subsection{Adaptive Runge--Kutta--Fehlberg Method}

The RKF45 method computes a fourth-order estimate $\widetilde y_{n+1}$ and a fifth-order estimate $y_{n+1}$ from six shared stage values.  Their difference estimates the local error:
\[
 \varepsilon_c=|y_{n+1}-\widetilde y_{n+1}|.
\]
If the current step is $h_c$ and the desired local tolerance is $\varepsilon_0$, a typical new step is
\[
 h_0=h_c\left(\frac{\varepsilon_0}{\varepsilon_c}\right)^{1/5}.
\]
If $\varepsilon_c\le\varepsilon_0$, the step is accepted and the next step size is adjusted.  Otherwise the result is rejected and recomputed with the smaller step.

\section{Multistep Methods}

Multistep methods use several previous function values.  Explicit Adams--Bashforth formulas predict a new value, while implicit Adams--Moulton formulas involve the unknown new value and are commonly used as correctors.

The four-step Adams--Bashforth formula is
\[
 y_{n+1}=y_n+\frac h{24}
 \left(55f_n-59f_{n-1}+37f_{n-2}-9f_{n-3}\right),
\]
where $f_j=f(x_j,y_j)$.

The three-step Adams--Moulton corrector of matching fourth order is
\[
 y_{n+1}=y_n+\frac h{24}
 \left(9f_{n+1}+19f_n-5f_{n-1}+f_{n-2}\right).
\]
A predictor--corrector implementation uses
\[
 y_{n+1}^{*}=y_n+\frac h{24}
 \left(55f_n-59f_{n-1}+37f_{n-2}-9f_{n-3}\right)
\]
and then
\[
 y_{n+1}=y_n+\frac h{24}
 \left(9f(x_{n+1},y_{n+1}^{*})+19f_n-5f_{n-1}+f_{n-2}\right).
\]
The method is not self-starting; values $y_1,y_2,y_3$ are normally generated with RK4.

\begin{example}
Use the fourth-order Adams--Bashforth/Adams--Moulton method with $h=0.2$ to approximate $y(0.8)$ for
\[
 y'=x+y,
 \qquad y(0)=1.
\]
\end{example}

\begin{solution}[6.2cm]
Using RK4 for the starting values gives
\[
 y_0=1,
 \quad y_1=1.242800,
 \quad y_2=1.583636,
 \quad y_3=2.044213.
\]
The corresponding function values are
\[
 f_0=1,
 \quad f_1=1.442800,
 \quad f_2=1.983636,
 \quad f_3=2.644213.
\]
The Adams--Bashforth predictor is
\[
 y_4^*=y_3+\frac{0.2}{24}
 (55f_3-59f_2+37f_1-9f_0)
 =2.650720.
\]
Then
\[
 f_4^*=0.8+y_4^*=3.450720,
\]
and the Adams--Moulton corrector gives
\[
 y_4=y_3+\frac{0.2}{24}
 (9f_4^*+19f_3-5f_2+f_1)
 =2.651056.
\]
Therefore
\[
 y(0.8)\approx2.6510.
\]
\end{solution}

\section{Higher-Order Initial-Value Problems}

A higher-order ODE can be rewritten as a first-order system.  For example,
\[
 y''=f(x,y,y'),
 \qquad y(x_0)=y_0,
 \qquad y'(x_0)=u_0,
\]
becomes
\[
 y'=u,
 \qquad
 u'=f(x,y,u).
\]
Any first-order numerical method can then be applied componentwise.

\subsection{Euler Method for a System}

For
\[
 y'=g(x,y,u),
 \qquad
 u'=f(x,y,u),
\]
Euler's method is
\[
 y_{n+1}=y_n+hg(x_n,y_n,u_n),
\]
\[
 u_{n+1}=u_n+hf(x_n,y_n,u_n).
\]

\begin{example}
Use Euler's method in two steps to approximate $y(0.2)$ and $y'(0.2)$ for
\[
 y''+xy'+y=0,
 \qquad y(0)=1,
 \qquad y'(0)=2.
\]
\end{example}

\begin{solution}[4.8cm]
Let $u=y'$.  Then
\[
 y'=u,
 \qquad
 u'=-xu-y.
\]
With $h=0.1$ and $(y_0,u_0)=(1,2)$,
\[
 y_1=1+0.1(2)=1.2,
 \qquad
 u_1=2+0.1[-0(2)-1]=1.9.
\]
At $x_1=0.1$,
\[
 y_2=1.2+0.1(1.9)=1.39,
\]
\[
 u_2=1.9+0.1[-0.1(1.9)-1.2]=1.761.
\]
Thus
\[
 y(0.2)\approx1.39,
 \qquad
 y'(0.2)\approx1.761.
\]
\end{solution}

\subsection{RK4 for a System}

For
\[
 y'=p(x,y,u),
 \qquad
 u'=q(x,y,u),
\]
RK4 computes paired stages $(k_i,m_i)$:
\[
 k_1=hp(x_n,y_n,u_n),
 \qquad
 m_1=hq(x_n,y_n,u_n),
\]
\[
 k_2=hp\left(x_n+\frac h2,y_n+\frac{k_1}{2},u_n+\frac{m_1}{2}\right),
\]
\[
 m_2=hq\left(x_n+\frac h2,y_n+\frac{k_1}{2},u_n+\frac{m_1}{2}\right),
\]
with analogous formulas for $(k_3,m_3)$ and $(k_4,m_4)$.  Then
\[
 y_{n+1}=y_n+\frac16(k_1+2k_2+2k_3+k_4),
\]
\[
 u_{n+1}=u_n+\frac16(m_1+2m_2+2m_3+m_4).
\]

\begin{example}
For
\[
 y''=4(y-x),
 \qquad y(0)=0,
 \qquad y'(0)=0,
\]
use one RK4 step with $h=1/3$ to approximate $y(1/3)$ and $y'(1/3)$.
\end{example}

\begin{solution}[5.0cm]
Let $u=y'$.  The system is
\[
 y'=u,
 \qquad
 u'=4(y-x).
\]
Applying the paired RK4 stages over $h=1/3$ gives
\[
 y_1\approx-0.02469,
 \qquad
 u_1\approx-0.23045.
\]
Therefore
\[
 y\left(\frac13\right)\approx-0.0247,
 \qquad
 y'\left(\frac13\right)\approx-0.2305.
\]
\end{solution}

\section{Stability and Stiffness}

A numerical method is stable if local perturbations and round-off errors do not grow catastrophically as the computation proceeds.

For the test equation
\[
 y'=-\lambda y,
 \qquad \lambda>0,
\]
Euler's method gives
\[
 y_{n+1}=(1-\lambda h)y_n.
\]
Stability requires
\[
 |1-\lambda h|\le1,
\]
so
\[
 0\le h\le\frac2\lambda.
\]
For a system
\[
 \vect y'=-A\vect y
\]
with positive eigenvalues, Euler stability requires approximately
\[
 h\le\frac{2}{\lambda_{\max}}.
\]

An IVP is stiff when some solution components decay or vary much more rapidly than others.  For
\[
 \vect y'=-A\vect y,
\]
widely separated eigenvalue magnitudes indicate stiffness.  An explicit method may then require an extremely small step for stability even when the slowly varying part of the solution would not require it for accuracy.

\begin{example}
For
\[
 y''+1001y'+1000y=0,
\]
find a step-size restriction for Euler's method to be stable.
\end{example}

\begin{solution}[3.8cm]
The characteristic polynomial is
\[
 r^2+1001r+1000=(r+1)(r+1000),
\]
so the decay rates are $1$ and $1000$.  The largest eigenvalue magnitude governing Euler stability is $1000$.  Therefore
\[
 h\le\frac2{1000}=0.002.
\]
The very different rates $1$ and $1000$ also show that the problem is stiff.
\end{solution}
